\documentclass[a4paper]{article}
\usepackage[warn]{mathtext}
\usepackage{graphicx} % Required for inserting images
\usepackage{wrapfig}
\usepackage{enumitem}
\usepackage[center]{caption}
\usepackage[english, russian]{babel}
\usepackage{amsmath}
\usepackage{geometry}
\setcounter{totalnumber}{5}

\geometry{verbose, a4paper, tmargin=2cm, bmargin=2cm, lmargin=2.0cm, rmargin=2.0cm}

\begin{document}
	
	\begin{titlepage}
		\centering
		{\scshape\Large МОСКОВСКИЙ ФИЗИКО-ТЕХНИЧЕСКИЙ ИНСТИТУТ \\
			(НАЦИОНАЛЬНЫЙ ИССЛЕДОВАТЕЛЬСКИЙ УНИВЕРСИТЕТ)\\ % название ВУЗа большим шрифтом
			Физтех-школа аэрокосмических технологий}
		
		\vspace{4cm} % отступ 4 см по вертикали
		{\LARGE Отчет о выполнении лабораторной работы 2.5.1}
		\vspace{1cm} % ещё 1 см
		
		{\huge\bf Измерение коэффициента поверхностного натяжения жидкости}
		
		\vspace{1cm} % ещё 1 см
		\vfill % заполнение по вертикали пустотой (LaTeX сам решит, сколько здесь отступить)
		
		\begin{flushright} % Этот текст будет с правого края страницы
			{\LARGE Губарев Никита, Б03-502}
		\end{flushright}
		\vfill 
		\today % Дата сборки документа
	\end{titlepage}
	\tableofcontents 
	\section{Аннотация}
	1) измерение температурной зависимости  коэффициента поверхностного натяжения дистиллированной воды с использованием известного коэффициента поверхностного натяжения спирта;  
	
	\noindent 2) определение полной поверхностной энергии  и теплоты, необходимой для изотермического образования единицы  поверхности жидкости  при различной температуре. 
	\section{Описание работы}

	\subsection{Теоретические сведения}
	
	Наличие свободной поверхности жидкости приводит к возникновению сил поверхностного натяжения, стремящихся сократить площадь этой поверхности. Коэффициент поверхностного натяжения \(\sigma\) определяется как сила, действующая на единицу длины контура, ограничивающего поверхность. Для искривлённой границы раздела двух сред (например, жидкость–газ) разность давлений по разные стороны поверхности описывается формулой Лапласа. Для сферической поверхности радиусом \(r\):
	\[
	\Delta P = P_{\text{внутри}} - P_{\text{снаружи}} = \frac{2\sigma}{r}.
	\tag{1}
	\]
	
	В методе Ребиндера (метод максимального давления пузырька) измеряют избыточное давление \(\Delta P\), необходимое для выдувания пузырька газа в жидкость через капилляр известного радиуса. В момент отрыва пузырька его радиус становится равным радиусу капилляра \(r\), а давление достигает максимума. Если использовать эталонную жидкость с известным \(\sigma_{\text{эт}}\), то из (1) можно определить радиус капилляра:
	\[
	r = \frac{2\sigma_{\text{эт}}}{\Delta P_{\text{эт}}}.
	\tag{2}
	\]
	Для исследуемой жидкости при тех же условиях коэффициент поверхностного натяжения вычисляется по формуле
	\[
	\sigma = \frac{r}{2}\,\Delta P = \sigma_{\text{эт}}\,\frac{\Delta P}{\Delta P_{\text{эт}}}.
	\tag{3}
	\]
	
	При измерении температурной зависимости \(\sigma\) необходимо учитывать влияние теплового расширения жидкости и возможное различие температур в разных частях установки. Для этого кончик капилляра погружают на постоянную глубину \(h\) (обычно почти до дна сосуда). В этом случае манометр регистрирует сумму лапласовского давления и гидростатического давления столба жидкости:
	\[
	P = \frac{2\sigma}{r} + \rho g h.
	\tag{4}
	\]
	Произведение \(\rho h\) при изменении температуры остаётся практически постоянным, так как масса столба жидкости над кончиком капилляра не меняется. Гидростатическую добавку определяют экспериментально: при фиксированной температуре измеряют давление \(P_1\) при касании капилляром поверхности (\(h=0\)) и давление \(P_2\) при полном погружении (\(h = h_0\)). Тогда
	\[
	\rho g h_0 = P_2 - P_1.
	\tag{5}
	\]
	
	Температурная зависимость коэффициента поверхностного натяжения для многих жидкостей вдали от критической точки близка к линейной. Полная поверхностная энергия единицы площади \(U/F\) и теплота изотермического образования единицы поверхности \(q\) связаны с \(\sigma(T)\) соотношениями Гиббса–Гельмгольца:
	\[
	\frac{U}{F} = \sigma - T \frac{d\sigma}{dT},
	\qquad
	q = -T \frac{d\sigma}{dT}.
	\tag{6}
	\]
	
	\subsection{Схема эксперимента}
	
	Схема установки для измерения коэффициента поверхностного натяжения методом Ребиндера приведена в приложении 1. Основными элементами являются:
	\begin{itemize}
		\item колба \(B\) с исследуемой жидкостью (дистиллированной водой), помещённая в водяную рубашку \(D\), через которую прогоняется вода из термостата для поддержания заданной температуры;
		\item колба \(E\) с эталонной жидкостью (этиловым спиртом), используемой для калибровки;
		\item металлическая игла \(C\), погружаемая в жидкость; верхний конец иглы открыт в атмосферу, нижний — в жидкость;
		\item аспиратор \(A\), создающий разряжение в системе; краны \(K_1\) и \(K_2\) позволяют регулировать поток воды и, следовательно, скорость снижения давления;
		\item спиртовой микроманометр для измерения разности давлений между атмосферой и разряженным объёмом.
	\end{itemize}
	
	
	Измерения проводят в герметичной системе. Аспиратор заполняют водой, после чего открывают кран \(K_1\); вода вытекает, создавая разряжение в нижней полости аспиратора и, через соединительные трубки, в колбах \(B\) и \(E\). По мере снижения давления пузырьки воздуха начинают проходить через жидкость в колбе, в которую погружена игла. Микроманометр измеряет максимальное давление в момент отрыва пузырька. Для получения воспроизводимых результатов частоту образования пузырьков выбирают не более одного пузырька за 5 секунд.
	
	Калибровка установки проводится по этиловому спирту, коэффициент поверхностного натяжения которого известен. Измерив \(\Delta P_{\text{спирт}}\) и зная \(\sigma_{\text{спирт}}\), по формуле (2) находят эффективный радиус иглы. Затем иглу переносят в колбу с водой и проводят измерения \(\sigma(T)\) в диапазоне температур от комнатной до \(60^{\circ}\text{C}\). Для исключения влияния теплового расширения и неравномерности нагрева иглу погружают на постоянную глубину; гидростатическую поправку определяют, как описано в разделе 2.1, сравнивая показания манометра при двух положениях иглы (на поверхности и на дне).
	\section{Методика измерений}
	
	\begin{enumerate}
		\item Проверьте герметичность установки. Для этого заполните аспиратор водой. Чистую сухую иглу установите в сосуд со спиртом так, чтобы кончик иглы лишь касался поверхности спирта. Плотно закройте обе колбы В и Е пробками. Откройте кран К1 аспиратора и добейтесь пробулькивания пузырьков воздуха в колбе. Замерьте показания микроманометра. Закройте кран К1. Наблюдайте за показаниями манометра: при отсутствии течи в установке столбик спирта в манометре будет неподвижен. При наличии течи следует обратиться к преподавателю или лаборанту.
		
		\item Убедившись в герметичности системы, начните измерения. Откройте кран К1. Подберите частоту падения капель из аспиратора так, чтобы максимальное давление манометра не зависело от этой частоты (не чаще, чем 1 капля в 5 секунд).
		
		\item Измерьте максимальное давление \(\Delta P_{\text{спирт}}\) при пробулькивании пузырьков воздуха через спирт. По разбросу результатов оцените случайную погрешность измерения. Пользуясь табличным значением коэффициента поверхностного натяжения спирта, определите по формуле (1) диаметр иглы. Сравните полученный результат с диаметром иглы, измеренным по микроскопу.
		
		\item Перенесите предварительно промытую и просушенную от спирта иглу в колбу с дистиллированной водой. Измерьте максимальное давление \(P_1\) при пробулькивании пузырьков, когда игла лишь касается поверхности воды. Аспиратор должен быть предварительно заполнен водой почти доверху. Отрегулируйте скорость поднятия уровня спирта в манометре и сохраняйте её в течение всех экспериментов. Измерьте расстояние между верхним концом иглы и любой неподвижной частью прибора \(h_1\).
		
		\item Утопите иглу до предела (между концом иглы и дном необходимо оставить небольшой зазор, чтобы образующийся пузырёк не касался дна). Измерьте \(h_2\) (как в пункте 4). Измерьте максимальное давление в пузырьках \(P_2\). По разности давлений \(\Delta P = P_2 - P_1\) определите глубину погружения \(\Delta h\) иглы и сравните с \(\Delta h = h_1 - h_2\).
		
		\item Снимите температурную зависимость \(\sigma(T)\) дистиллированной воды. Для этого включите термостат и подождите, пока нужная вам температура не стабилизируется. Следует заметить, что термометр показывает температуру воды в термостате. Время установления заданной температуры в колбе с исследуемой жидкостью довольно большое, поэтому после установления заданной температуры в термостате дождитесь прогрева исследуемой жидкости в колбе в течение нескольких минут (5–7 минут). После этого проведите измерение давления. Для уменьшения погрешности опыта замер давления при фиксированной температуре следует проводить несколько раз.
		
		Проводить измерение температурной зависимости рекомендуется в диапазоне \(20^{\circ}\text{C}\) – \(60^{\circ}\text{C}\) через \(5^{\circ}\). Запрещается нагревать воду в термостате выше \(60^{\circ}\text{C}\).
		
		\item Оцените погрешность измерения давления и температуры. Рассчитайте величину коэффициента поверхностного натяжения воды \(\sigma(T)\), используя значение диаметра иглы, полученное при измерениях на спирте (или измеренное на микроскопе).
		
		\item Постройте график зависимости \(\sigma(T)\) и определите по графику температурный коэффициент \(\frac{d\sigma}{dT}\). Оцените точность результата.
		
		\item На другом графике постройте зависимость от температуры:
		\begin{enumerate}
			\item теплоты образования единицы поверхности жидкости \(q = -T \cdot \frac{d\sigma}{dT}\);
			\item поверхностной энергии \(U\) единицы площади \(F\): \(\frac{U}{F} = \sigma - T \cdot \frac{d\sigma}{dT}\).
		\end{enumerate}
	\end{enumerate}

	
	\section{Обработка результатов}
	
	\subsection{Определение радиуса иглы по эталонной жидкости}
	
	Для калибровки используется этиловый спирт с табличным значением коэффициента поверхностного натяжения
	\(\sigma_{\text{спирт}} = 22.3\) мН/м при комнатной температуре.
	Измеренное максимальное давление при пробулькивании пузырьков через спирт составляет
	\(\Delta P_{\text{спирт}} = 41 \pm 1\) отсчёт. Цена деления микроманометра
	\(K = 0.2 \cdot 9.81 = 1.962\) Па/отсчёт, поэтому
	\[
	\Delta P_{\text{спирт}} = 41 \cdot 1.962 = 80.44 \pm 1.96\ \text{Па}.
	\]
	
	Из формулы Лапласа \(\Delta P = \dfrac{2\sigma}{r}\) выражаем радиус:
	\[
	r = \frac{2\sigma_{\text{спирт}}}{\Delta P_{\text{спирт}}}.
	\]
	Подставляя численные значения:
	\[
	r = \frac{2 \cdot 0.0223}{80.44} = 5.545 \cdot 10^{-4}\ \text{м} = 0.5545 \pm 0.0135\ \text{мм}.
	\]
	
	
	\begin{table}[h!]
		\centering
		\caption{Калибровка радиуса иглы по эталонной жидкости}
		\label{tab:calib}
		\begin{tabular}{|c|c|c|}
			\hline
			Параметр & Значение & Погрешность\\
			\hline
			\(\Delta P_{\text{спирт}}\) & 80.44 & 1.96 Па \\
			\(\sigma_{\text{спирт}}\) & 22.3 мН/м &\\
			\(r\) & 0.554 & 0.014  мм \\
			\(d\) & 1.108 & 0.028  мм \\
			\hline
		\end{tabular}
	\end{table}
	Измеренный по микроскопу диаметр иглы: \(d_{\text{микро}} = 1.10 \pm 0.05\) мм.
	\newpage
	\subsection{Определение гидростатической добавки}
	
	При температуре \(t = 23.5^\circ\text{C}\) проведены два измерения (табл.~\ref{tab:hydro_meas}).
	
	\begin{table}[h!]
		\centering
		\caption{Измерения для определения гидростатической поправки}
		\label{tab:hydro_meas}
		\begin{tabular}{|c|c|c|c|}
			\hline
			Положение иглы & Давление, Па & Высота \(h\), мм \\
			\hline
			Касание поверхности & \(223.67 \pm 1.96\) & \(19 \pm 1\) \\
			Погружение на дно & \(339.43 \pm 1.96\) & \(7 \pm 1\) \\
			\hline
		\end{tabular}
	\end{table}
	
	Глубина погружения:
	\[
	\Delta h = h_1 - h_2 = 12 \pm 2\ \text{мм}.
	\]
	
	Гидростатическое давление:
	\[
	P_{\text{гидр}} = P_2 - P_1 = 115.76 \pm 3.92\ \text{Па}.
	\]
	
	Теоретическое значение:
	\[
	P_{\text{гидр, теор}} = \rho g \Delta h = 1000 \cdot 9.81 \cdot 0.012 = 117.72 \pm 13.73\ \text{Па}.
	\]
	Расхождение находится в пределах погрешности. В дальнейшем используем экспериментальное значение
	\(P_{\text{гидр}} = 115.76\ \text{Па}\) с погрешностью \(2.77\ \text{Па}\).
	
	\subsection{Расчёт коэффициента поверхностного натяжения воды}
	
	Для каждой температуры измерено давление \(P_2\) (игла погружена на ту же глубину).
	Лапласовское давление (вклад искривлённой поверхности) находим как
	\[
	P_{\text{лапл}} = P_2 - P_{\text{гидр}}.
	\]
	При \(t = 23.5^\circ\text{C}\) используем \(P_1\), так как в этом случае игла касается поверхности
	и гидростатическая добавка отсутствует.
	
	Коэффициент поверхностного натяжения воды вычисляем через радиус иглы:
	\[
	\sigma = \frac{P_{\text{лапл}} \cdot r}{2}.
	\]

	
	\begin{table}[h!]
		\centering
		\caption{Коэффициент поверхностного натяжения воды при разных температурах}
		\label{tab:sigma_water}
		\begin{tabular}{|c|c|c|c|c|c|}
			\hline
			\(t\), \(^\circ\)C & \(P_2\), Па & \(P_{\text{лапл}}\), Па & \(\Delta P_{\text{лапл}}\), Па & \(\sigma\), мН/м & \(\Delta\sigma\), мН/м \\
			\hline
			23.5 & \(223.67\)* & \(223.67\) & \(1.96\) & \(62.00\) & \(1.6\) \\
			30   & \(319.81\)  & \(204.05\) & \(3.40\) & \(56.57\) & \(1.7\) \\
			40   & \(311.96\)  & \(196.20\) & \(3.40\) & \(54.39\) & \(1.6\) \\
			50   & \(306.07\)  & \(190.31\) & \(3.40\) & \(52.76\) & \(1.6\) \\
			60   & \(300.19\)  & \(184.43\) & \(3.40\) & \(51.13\) & \(1.6\) \\
			\hline
			\multicolumn{6}{l}{*При \(23.5^\circ\text{C}\) указано \(P_2 = P_1\), \(P_{\text{лапл}} = P_1\).} \\
		\end{tabular}
	\end{table}
	
	\subsection{Температурная зависимость \(\sigma(T)\) и её производная}
	
	Для определения температурного коэффициента \(\frac{d\sigma}{dT}\) проведена линейная аппроксимация
	зависимости \(\sigma(T)\) методом взвешенных наименьших квадратов (веса \(w_i = 1/\Delta\sigma_i^2\)).
	Аппроксимирующая функция: \(\sigma(T) = a + bT\).
	В результате получено:
	\[
	b = \frac{d\sigma}{dT} = -0.30 \pm 0.02\ \text{мН/(м·К)}.
	\]
	

	
	\subsection{Теплота образования единицы поверхности и полная поверхностная энергия}
	
	Теплота, необходимая для изотермического образования единицы поверхности:
	\[
	q = -T \frac{d\sigma}{dT}.
	\]
	Полная поверхностная энергия единицы площади (внутренняя энергия поверхности):
	\[
	\frac{U}{F} = \sigma - T \frac{d\sigma}{dT}.
	\]
	
	Для расчёта использовано значение производной \(b = -0.300\ \text{мН/(м·К)}\).
	Погрешности величин вычислены по формулам:
	\[
	\Delta q = T \cdot \Delta b,\qquad
	\Delta(U/F) = \sqrt{(\Delta\sigma)^2 + (T \cdot \Delta b)^2}.
	\]
	
	Результаты приведены в табл.~\ref{tab:thermo}.
	
	\begin{table}[h!]
		\centering
		\caption{Термодинамические характеристики поверхности воды}
		\label{tab:thermo}
		\begin{tabular}{|c|c|c|c|c|c|c|}
			\hline
			\(t\), \(^\circ\)C & \(T\), К & \(\sigma\), мН/м & \(q\), мН/м & \(\Delta q\), мН/м & \(U/F\), мН/м & \(\Delta(U/F)\), мН/м \\
			\hline
			23.5 & 296.65 & 62.00 & 89.0 & 5.9 & 151.0 & 6.1 \\
			30   & 303.15 & 56.57 & 90.9 & 6.1 & 147.5 & 6.3 \\
			40   & 313.15 & 54.39 & 93.9 & 6.3 & 148.3 & 6.5 \\
			50   & 323.15 & 52.76 & 96.9 & 6.5 & 149.7 & 6.6 \\
			60   & 333.15 & 51.13 & 99.9 & 6.7 & 151.0 & 6.8 \\
			\hline
		\end{tabular}
	\end{table}
	
	
	\section{Выводы}
		Полученные значения коэффициента поверхностного натяжения воды ниже табличных
	(при \(20^\circ\text{C}\) \(\sigma = 72.75\) мН/м, при \(60^\circ\text{C}\) \(\sigma \approx 66.2\) мН/м).
	Такое систематическое расхождение может быть связано с:
	\begin{itemize}
		\item использованием спирта в качестве эталонной жидкости без учёта возможных примесей;
		\item наличием поверхностно-активных веществ в дистиллированной воде;
		\item отличием эффективного радиуса пузырька от геометрического радиуса иглы из-за смачивания.
	\end{itemize}
	Тем не менее, линейный характер зависимости \(\sigma(T)\) и рассчитанные величины \(q\) и \(U/F\)
	качественно соответствуют теории: с ростом температуры \(\sigma\) убывает, теплота образования
	поверхности возрастает, а полная поверхностная энергия остаётся практически постоянной,
	что согласуется с представлениями о слабой температурной зависимости внутренней энергии
	поверхностного слоя.
	\newpage
	\section{Приложения}
	\begin{figure}[h]
		\centering
		\includegraphics[width=0.8\textwidth]{1.png}
		\caption*{Приложение 1. Схема установки}
	\end{figure}
	
	\begin{figure}[h]
		\centering
		\includegraphics[width=0.9\textwidth]{2.png}
		\caption*{Приложение 2. Графики зависимостей поверхностного натяжения, теплоты образования единицы поверхности жидкости, поверхностной энергии от температуры}
	\end{figure}
	

\end{document}